\documentclass[11pt]{article}
\usepackage[margin=1in]{geometry}
\usepackage[T1]{fontenc}
\usepackage{lmodern}
\usepackage{microtype}
\usepackage{booktabs}
\usepackage{tabularx}
\usepackage{array}

\begin{document}

\section{Results and Discussion}
\label{sec:results-discussion}

The results below are taken from \texttt{results\_hybrid/summary.txt} and
\texttt{results\_hybrid/comparison.csv}. The principal comparison uses 100
jobs and 20 AGVs on each of ten job sets (seeds 1--10). Configuration A is
the centralized paper-style planner, B is centralized planning with lift
lobbies, and C is the hybrid central-plus-floor-server planner using its
default first-come-first-served leg order. The percentage column reports C
relative to A; a negative value indicates a reduction.

\subsection{Correctness and decomposition check}

All evaluated plans in configurations A, B, and C have \texttt{ok=True} in
\texttt{comparison.csv}; the saved summary records that every plan passes the
independent \texttt{agv\_core.verify()} checker. The separate decomposition
check runs the hybrid planner in priority order and compares it with
configuration B. The plans match on every one of the ten seeds. This supports
that, under priority-order processing and the tested model, splitting routes
between floor servers preserves the centralized-with-lobbies schedule. It
does not imply equality for configuration C's default FCFS order, which is the
order used for the hybrid performance results below.

\subsection{Mean performance over ten seeds}

Table~\ref{tab:mean-results} gives the saved mean values. In the tested setup,
configuration C reduces last-delivery makespan by 65\%, total tardiness by
72\%, usage-time imbalance by 80\%, total waiting by 58\%, and longest
single-AGV waiting by 74\% relative to A. The busiest-server CPU time falls
from 77.2~ms to 5.3~ms (93\%); total CPU summed over all servers falls by
78\%, from 77.2~ms to 17.1~ms. B also improves on A, indicating that lobby
waiting points contribute to the observed results. Comparing C with B adds
the effects of floor-server decomposition and FCFS leg processing.

\begin{table}[htbp]
\centering
\small
\caption{Mean metrics over 10 seeds; C vs. A percentages are as reported in
the saved summary.}
\label{tab:mean-results}
\begin{tabularx}{\textwidth}{>{\raggedright\arraybackslash}Xrrrr}
\toprule
Metric & A & B & C & C vs. A \\
\midrule
Last-delivery makespan (s) & 7,972.4 & 3,705.9 & 2,768.9 & $-65\%$ \\
Total tardiness (s) & 316,893.6 & 122,278.8 & 90,104.5 & $-72\%$ \\
Load imbalance (s) & 7,149.9 & 2,832.3 & 1,441.2 & $-80\%$ \\
Total waiting (s) & 64,689.7 & 27,647.7 & 27,104.5 & $-58\%$ \\
Longest AGV wait (s) & 7,133.4 & 2,830.8 & 1,887.7 & $-74\%$ \\
Busiest-server CPU (ms) & 77.2 & 15.5 & 5.3 & $-93\%$ \\
Total server CPU (ms) & 77.2 & 15.5 & 17.1 & $-78\%$ \\
Windows on busiest server & 2,030.7 & 1,733.5 & 634.5 & $-69\%$ \\
Central--floor messages & 0.0 & 0.0 & 337.2 & --- \\
Algorithm 3 calls & 106.6 & 165.1 & 168.6 & $+58\%$ \\
\bottomrule
\end{tabularx}
\end{table}

The CSV additionally records a mean of 2,095.8 transmitted window records
(\texttt{payload\_windows}) for C, compared with zero for A and B, which do
not exchange central--floor messages. C therefore introduces explicit
coordination traffic: on average 337.2 messages and about 2,096 window
records per 100-job run. These are measured message and payload counts, not
network bytes or communication latency.

\subsection{Fleet-size sweep}

Table~\ref{tab:fleet-sweep} reports mean load imbalance and total tardiness
over the same seeds as fleet size increases. In all three configurations,
mean tardiness falls as more AGVs are available. For A, imbalance also
declines over the sweep. In B and C, imbalance increases with fleet size,
although C remains lowest at every tested fleet size. The increase reflects
the metric's range across AGV return times: adding vehicles does not
necessarily distribute the greedy job assignments or conflict-induced waits
evenly. At 20 AGVs, C has mean imbalance of 1,441~s and tardiness of
90,104~s, compared with 7,150~s and 316,894~s for A.

\begin{table}[htbp]
\centering
\small
\caption{Fleet-size sweep: mean load imbalance / total tardiness over 10 seeds.}
\label{tab:fleet-sweep}
\begin{tabular}{rrrr}
\toprule
AGVs & A (s / s) & B (s / s) & C (s / s) \\
\midrule
4  & 9,119 / 602,120 & 2,043 / 256,490 & 518 / 235,718 \\
8  & 8,394 / 453,005 & 2,346 / 149,864 & 788 / 126,223 \\
12 & 8,249 / 421,464 & 2,456 / 125,292 & 1,145 / 100,319 \\
16 & 7,466 / 348,159 & 2,715 / 117,682 & 1,230 / 90,343 \\
20 & 7,150 / 316,894 & 2,832 / 122,279 & 1,441 / 90,104 \\
\bottomrule
\end{tabular}
\end{table}

\subsection{Per-server state and computation}

For seed~7, the saved server breakdown is shown in Table~\ref{tab:seed7}.
The central server in C holds 146 windows and uses 5.0~ms CPU; the three
floor servers hold 478, 619, and 454 windows and use 2.9, 5.6, and 4.3~ms,
respectively. Thus the largest individual reservation set is 619 windows,
compared with 1,989 windows on A's single central server. The maximum server
CPU is 5.6~ms for C, against 79.9~ms for A. The hybrid distributes the
per-floor node and track windows instead of requiring one planner to retain
the complete set. The sum of server CPU is not the same as the peak: C's
total is higher than B's mean CPU in Table~\ref{tab:mean-results}, while its
busiest-server time is substantially lower.

\begin{table}[htbp]
\centering
\small
\caption{Per-server CPU and reservation-window counts for seed 7.}
\label{tab:seed7}
\begin{tabular}{lrr}
\toprule
Server & CPU (ms) & Windows held \\
\midrule
A: central & 79.9 & 1,989 \\
C: central & 5.0 & 146 \\
C: floor 1 & 2.9 & 478 \\
C: floor 2 & 5.6 & 619 \\
C: floor 3 & 4.3 & 454 \\
\bottomrule
\end{tabular}
\end{table}

\subsection{Discussion}

The performance pattern is consistent with two changes in C: waits can be
absorbed at lift lobbies near the floor where a conflict occurs, and floor
servers resolve local node and track reservations within separate state
sets. A can only wait at the single parking node, while B adds lobbies but
still maintains centralized planning. The substantial A-to-B improvement
therefore shows that additional waiting locations matter in this factory;
the further improvement from B to C is associated with the hybrid's
decomposed scheduling and FCFS leg order. Since these changes interact, the
three configurations do not isolate every causal effect independently.

The state and CPU results indicate that decomposition reduces the largest
per-server planning burden for the tested three-floor, 20-AGV cases. It does
not eliminate coordination: C averages 337.2 central--floor messages per
run, with about 2,095.8 reservation-window records in the payload. Algorithm~3
is called 168.6 times on average in C, slightly more than B's 165.1 calls and
58\% more than A's 106.6. A plausible explanation is that FCFS leg planning
creates local scheduling decisions at floor servers and includes lobby
reservations, each of which can trigger a separate conflict-resolution call.
The extra calls are small relative to the reduction in peak CPU and stored
windows; the counts alone do not establish a causal explanation.

These results support the hybrid architecture for this lift-connected
factory model, particularly when lift contention and floor-local conflicts
contribute substantially to waiting. The fleet sweep shows lower tardiness
and imbalance for C than A and B at each tested fleet size, but does not
establish asymptotic scaling beyond 20 AGVs or three floors. The method also
does not remove limitations inherited from the baseline: job assignment is
greedy and priority-dependent, routes use shortest paths, lift trips are
booked under the model's fixed travel-time assumptions, and no global
optimality guarantee is provided. Communication and aggregate CPU remain
costs of decomposition, and the experiments do not report network latency or
validate physical multi-floor hardware. These limits motivate broader tests
with more floors, varied lift capacities and travel times, and different
traffic densities.

\end{document}
